Our starting point will be \autoref{eq:por-conv-work}:
\begin{align}
  \text{ConvWork}_{R\rightarrow L}(w) = \; & \frac{L_d}{L_r} \cdot X_{R\rightarrow L} \cdot \frac{R_r}{R_d} \cdot w
    & &\text{R-hashes }\rightarrow\text{ L-hashes}
    \nonumber
    % \\[0.5em]
%   = \; & \frac{L_d}{L_r} \cdot X_{R\rightarrow L} \cdot \frac{R_r}{R_d} \cdot w
%     & &\text{R-hashes }\rightarrow\text{ L-hashes}
\end{align}

It's easy to see that we can work with this -- of course, the \textsc{WeightOf} function needs access to R's difficulty and block reward, and the exchange rate, too.
Do we \emph{need} to use $R_d$, though?
Surely \emph{any} conversion to L-hashes/L-block will work, right?

Let's see:
\begin{align}
    & \frac{L_d}{L_r}
        & & \frac{\text{L-hashes}}{\text{L-coin}}
        \nonumber
        \\[0.5em]
    & \frac{L_d}{L_r} \cdot X_{R\rightarrow L}
        & & \frac{\text{L-hashes}}{\text{R-coin}}
        \nonumber
        \\[0.5em]
    & \frac{L_d}{L_r} \cdot X_{R\rightarrow L} \cdot R_r
        & & \frac{\text{L-hashes}}{\text{R-block}}
        \nonumber
        \\[0.5em]
    & \frac{L_d}{L_r} \cdot X_{R\rightarrow L} \cdot R_r \cdot \frac{R_f}{L_f}
        & & \frac{\text{L-hashes}}{\text{L-block}}
        \nonumber
        \\[0.5em]
    \therefore \text{ConvDEX}_{R\rightarrow L}(R_r) = \; & \frac{L_d}{L_r} \cdot X_{R\rightarrow L} \cdot R_r \cdot \frac{R_f}{L_f}
          & &\frac{\text{R-coins}}{\text{R-block}}\rightarrow\frac{\text{L-hashes}}{\text{L-block}}
          \label{eq:por-conv-dex}
\end{align}

Okay, so we \emph{can} convert to L's block-weight \emph{without knowing R's difficulty}.
It's good to know that we have \emph{options} there -- there's \emph{excess capacity} in the values we can convert between.
If one value is not available, then we can substitute a combination of other values.
(This is relevant for \autoref{sec:converting-confirmations}.)

\todoDraftOnly{This (above) is one of the properties that we should quickcheck}

To move forward, we need to use $X_{R\rightarrow L}$ in a new \textsc{WeightOf} function -- \emph{where does that value come from?}

The value can't be hard-coded, or provided by a third party -- these would introduce vulnerabilities.
We need an accurate exchange rate for the moment of conversion, too -- one that is reactive enough to remain up-to-date, and precise enough to be useful.
It also needs to be \emph{provable}; in practice, it needs to be \emph{on-chain} and available to both chains (and they should agree on the exchange rate).
If that wasn't the case, how could one chain validate another chain's PoRs?
Ideally, the exchange rate should \emph{already} be recorded in \emph{both} chains.
(This way, full nodes of either chain do not need extra data to calculate/verify chain-weight.)

\textbf{However,} there is a major loop-hole in some of the above requirements: they only matter for \emph{mutual PoR}.
What's different if we consider \emph{one-way PoR} instead?

In the case of one-way PoR, the reflecting chain doesn't need to know much about the reflected chain, and the heavy lifting (like the burden of knowing exchange rates) can be offloaded to \emph{the reflected chain only}.
Provided that the source of the exchange rate is a \emph{well-built}, protocol-level DEX, are there any real issues?

Let's consider this limited case first: one-way PoR where the \emph{reflected} chain hosts a DEX between it's root token, and the RT of the \emph{reflecting} chain.
Using \autoref{eq:por-conv-dex}, we can write \autoref{alg:weightof-dex-one-way}.

\begin{algorithm}
\caption{A \textsc{WeightOf} function for one-way PoR between chains with different root tokens}
\label{alg:weightof-dex-one-way}
\begin{algorithmic}
\Procedure{WeightOf}{$B_i$, $state$} \Comment{Convert reflecting weight to local via a DEX}
    \State $t \gets$ \Call{BlockTimestamp}{$B_i$}
    \State $C \gets$ \Call{ChainOf}{$B_i$, $state$}
    \State $C_r \gets$ \Call{BlockRewardOfChainAt}{$C$, $t$, $state$} \Comment{Block reward of $C$ at $t$}
    \State $C_f \gets$ \Call{BlockFrequencyOfChainAt}{$C$, $t$, $state$}
    \State $L_d \gets$ \Call{DifficultyOfLocalAt}{$t$, $state$}
    \State $L_f \gets$ \Call{BlockFrequencyOfLocalAt}{$t$, $state$}
    \State $L_r \gets$ \Call{BlockRewardOfLocalAt}{$t$, $state$}
    \State $X_{C\rightarrow L} \gets$ \Call{GetLocalDexRateFromAt}{$C$, $t$, $state$}
    % \State $w \gets$ \Call{$\text{ConvDEX}_{R\rightarrow L}$}{$L_d$, $L_f$, $L_r$, $C_f$, $C_r$, $X_{C\rightarrow L}$}
    \State \Return{\Call{$\text{ConvDEX}_{R\rightarrow L}$}{$L_d$, $L_f$, $L_r$, $C_f$, $C_r$, $X_{C\rightarrow L}$}}
\EndProcedure
\\
\Procedure{$\text{ConvDEX}_{R\rightarrow L}$}{$L_d$, $L_f$, $L_r$, $R_f$, $R_r$, $X_{R\rightarrow L}$} \Comment{Returns L-hashes/L-block}
    \State \Return{$(L_d \div L_r) \cdot X_{R\rightarrow L} \cdot R_r \cdot (R_f \div L_f)$} \Comment{See \autoref{eq:por-conv-dex}}
\EndProcedure
\end{algorithmic}
\end{algorithm}


\begin{algorithm}
\caption{A \textsc{WeightOf} function for PoR between chains with a different root token and equal block frequencies}
\label{alg:weightof-dex}
\begin{algorithmic}
\Procedure{WeightOf}{$B_i, state$} \Comment{The weight of a block normalized to coins}
    \State $t \gets$ \Call{BlockTimestamp}{$B_i$}
    \State $C \gets$ \Call{ChainOf}{$B_i$, $state$}
    \State $C_r \gets$ \Call{BlockRewardOfChainAt}{$C$, $t$, $state$} \Comment{Block reward of $C$ at $t$}
    % \State $C_f \gets$ \Call{BlockFrequencyOfChain}{$C$} \Comment{Block production Hz of $C$}
    \State $X_{C\rightarrow L} \gets$ \Call{ExchangeRateOfCoinToLocalAt}{$C$, $t$, $state$}
    \State \Return{$C_r \cdot X_{C\rightarrow L}$}
\EndProcedure
\end{algorithmic}
\end{algorithm}


% \begin{algorithm}
% \caption{A \textsc{WeightOf} function for one-way PoR between chains with different root tokens}
% \label{alg:weightof-dex-one-way}
% \begin{algorithmic}
% \Procedure{WeightOf}{$B_i$, $state$} \Comment{Convert reflecting weight to local via a DEX}
%     \State $t \gets$ \Call{BlockTimestamp}{$B_i$}
%     \State $C \gets$ \Call{ChainOf}{$B_i$, $state$}
%     \State $C_r \gets$ \Call{BlockRewardOfChainAt}{$C$, $t$, $state$} \Comment{Block reward of $C$ at $t$}
%     \State $C_f \gets$ \Call{BlockFrequencyOfChainAt}{$C$, $t$, $state$}
%     \State $L_d \gets$ \Call{DifficultyOfLocalAt}{$t$, $state$}
%     \State $L_f \gets$ \Call{BlockFrequencyOfLocalAt}{$t$, $state$}
%     \State $L_r \gets$ \Call{BlockRewardOfLocalAt}{$t$, $state$}
%     \State $X_{C\rightarrow L} \gets$ \Call{ExchangeRateOfCoinToLocalAt}{$C$, $t$, $state$}
%     % \State $w \gets$ \Call{$\text{ConvDEX}_{R\rightarrow L}$}{$L_d$, $L_f$, $L_r$, $C_f$, $C_r$, $X_{C\rightarrow L}$}
%     \State \Return{\Call{$\text{ConvDEX}_{R\rightarrow L}$}{$L_d$, $L_f$, $L_r$, $C_f$, $C_r$, $X_{C\rightarrow L}$}}
% \EndProcedure
% \\
% \Procedure{$\text{ConvDEX}_{R\rightarrow L}$}{$L_d$, $L_f$, $L_r$, $R_f$, $R_r$, $X_{R\rightarrow L}$} \Comment{Returns L-hashes/L-block}
%     \State \Return{$(L_d \div L_r) \cdot X_{R\rightarrow L} \cdot R_r \cdot (R_f \div L_f)$} \Comment{See \autoref{eq:por-conv-dex}}
% \EndProcedure
% \end{algorithmic}
% \end{algorithm}

This \emph{seems} okay; of course, it \emph{still depends on the DEX}.
That means that an attack on the DEX might be a way to attack the consensus algorithm.
Did we just introduce a vulnerability?

If an attacker \emph{reduces} the exchange rate (L-coins/R-coin) somehow, then that will \emph{decrease} the weight of reflected blocks.
For example, an attacker might sell \emph{a lot} of L-coins all at once (or cause that to happen through traditional market manipulation).
If the exchange rate goes down, then the \emph{rate of work} (i.e., how fast chain-weight is accumulated) of Chain L will \emph{seem} to decrease.

However, this isn't \emph{always} relevant: if Chain R is \emph{much} more secure than Chain L, then this doesn't change much in practice.
Say that Chain R contributes $\sim 100\times$ as much chain-weight as L does.
For an attacker targeting Chain L, reducing the exchange rate \emph{by half} means that the weight of reflecting blocks is halved.
So the attacker would need to add $\sim 50\times$ the \emph{current} hash-rate of L to compete against the reflections from R.
If the attacker was mining in \emph{public}, that would increase the local difficulty by $\sim 50\times$, which \emph{increases} the weight of reflecting blocks by $\sim 50\times$, too.
Of course, the attacker \emph{could} still mine in private,\footnote{
    For the purpose of UT and Amaroo, this turns out to be a non-issue.
    See \autoref{sec:the-simplex}, \autoref{sec:dapp-chain-security}, \autoref{sec:dos-and-dags}, and \autoref{sec:miner-resonance}.
} provided the exchange rate stays low.
Even though the attacker gets a $50\%$ discount via market manipulation, this is still a very difficult attack (and very noticeable).
In effect, attacking L is at least as difficult as attacking R.

If an attacker \emph{increases} the exchange rate (L-coins/R-coin) somehow, then that will \emph{increase} the weight of reflecting blocks.
This does not help the attacker.

It seems that \emph{in principle} there are some cases where it's safe to use $\text{ConvDEX}_{R\rightarrow L}$ for one-way PoR.
We'll revisit this later in \autoref{sec:dapp-chain-security}.

What about mutual PoR, though?
Notice that \autoref{alg:weightof-dex-one-way} is trivially generalized for mutual reflection by replacing \textsc{GetLocalDexRateFromAt} with a suitable \textsc{ExchangeRateToLocalOfAt}, i.e., a DEX that satisfies the requirements we specified earlier.
So it seems like our \emph{conversion method} should work; the weak-link is the DEX, not the conversion.

In the context of \emph{Ultra Terminum} and \emph{Amaroo}, this isn't an important problem to solve.
After all, this is a paper about \emph{consensus}, not trustless, decentralized, cross-chain, protocol-level markets.
If \emph{mutual PoR} is ever used to secure multiple chains with heterogenous tokens, this problem will need to be answered, though.
